% !TEX root = ../../main.tex
% C3.1 — Bài toán tìm kiếm người và truy hồi đa phương thức (RÚT GỌN 2026-09-29; bản cũ ở report/archive/)
% Nội dung: ReID (ảnh→ảnh) và text-based person search (văn bản→ảnh); tìm theo thuộc tính; khác nhận dạng khuôn mặt;
% hướng tiếp cận của đề tài (đơn vị kết quả là track, một không gian biểu diễn chung cho ba hình thức mô tả).
\section{Bài toán tìm kiếm người và truy hồi đa phương thức}
\label{sec:3.1}

\subsection{Phát biểu bài toán}
\label{sec:3.1.1}

Cho một tập hình ảnh người được trích xuất từ video của nhiều camera, gọi là \emph{tập dữ liệu tìm kiếm} (gallery) $\mathcal{G} = \{g_1, g_2, \ldots, g_n\}$, và một \emph{truy vấn} $q$ mô tả người cần tìm, hệ thống cần xếp hạng các phần tử của $\mathcal{G}$ theo mức độ phù hợp với $q$ và trả về $k$ phần tử phù hợp nhất. Cách tiếp cận phổ biến là học một hàm biểu diễn ánh xạ truy vấn và từng phần tử vào một không gian vector, rồi đo mức độ phù hợp bằng một hàm tương đồng $s(q, g_i)$; việc tìm kiếm khi đó trở thành tìm các vector gần nhất.

Trong video giám sát, người không có sẵn dưới dạng ảnh cắt gọn như các bộ dữ liệu chuẩn mà nằm trong khung hình toàn cảnh cùng nhiều người khác, nên bài toán còn gồm bước phát hiện. Xiao và cộng sự học đồng thời phát hiện và nhận dạng~\cite{xiao2017personsearch}, còn một hướng khác tách bộ phát hiện khỏi mô hình nhận dạng~\cite{zheng2017reidwild}; đề tài theo hướng sau để thay thế độc lập từng thành phần.

\subsection{Ba hình thức truy vấn}
\label{sec:3.1.2}
\label{sec:3.1.3}
\label{sec:3.1.4}

\paragraph{Tìm bằng hình ảnh.} Khi truy vấn là một ảnh chứa người cần tìm, bài toán tương ứng với \emph{tái nhận dạng người} (Person Re-identification, ReID): xác định các ảnh thuộc cùng một người trong dữ liệu thu từ nhiều camera hoặc nhiều thời điểm~\cite{ye2022reidsurvey}. Thách thức chính là sự thay đổi góc nhìn, ánh sáng, tư thế, hiện tượng che khuất và độ phân giải thấp: cùng một người có thể trông rất khác giữa hai camera, trong khi hai người mặc giống nhau lại có thể trông rất giống nhau.

\paragraph{Tìm bằng mô tả văn bản.} Khi người dùng chỉ có lời mô tả, ví dụ ``a man wearing a black jacket and blue jeans, carrying a backpack'', bài toán là tìm kiếm người bằng văn bản (Text-Based Person Search, TBPS), được Li và cộng sự đề xuất cùng bộ dữ liệu CUHK-PEDES gồm ảnh người kèm mô tả tiếng Anh~\cite{li2017cuhkpedes}. Đây là bài toán liên phương thức: mô hình phải đưa văn bản và ảnh về cùng một không gian biểu diễn, trong khi mô tả thường mơ hồ và khác biệt giữa những người thường nằm ở chi tiết nhỏ.

\paragraph{Tìm theo thuộc tính.} Người dùng chọn thuộc tính như giới tính, loại và màu trang phục, vật mang theo. Có thể dùng một mô hình nhận dạng thuộc tính người đi bộ~\cite{wang2022parsurvey} để lọc theo nhãn, cần dữ liệu gán nhãn cho từng thuộc tính, hoặc ghép thuộc tính thành câu theo mẫu cố định rồi tìm như văn bản. Đề tài chọn cách sau để dùng chung mô hình văn bản, nên thuộc tính ảnh hưởng thứ hạng chứ không phải điều kiện lọc bắt buộc. Lựa chọn phủ định như ``không mang ba lô'' không được hỗ trợ, vì các mô hình ngôn ngữ-thị giác kiểu CLIP gặp khó khăn đáng kể với phủ định, nhiều trường hợp chỉ đạt mức ngẫu nhiên~\cite{alhamoud2025negation}.

Khác nhận dạng khuôn mặt, tìm theo ngoại hình không dựa trên sinh trắc học, vốn cần khuôn mặt rõ, mà dựa trên trang phục, màu sắc và vật dụng, dễ thấy từ camera cao và xa nhưng có thể thay đổi và trùng giữa nhiều người. Kết quả vì vậy là gợi ý để đánh giá, không phải kết luận về danh tính.

\subsection{Hướng tiếp cận của đề tài}
\label{sec:3.1.6}

Bảng~\ref{tab:query-modes} tóm tắt ba hình thức truy vấn và cách đề tài xử lý.

\begin{table}[H]
\centering
\caption{Ba hình thức mô tả người cần tìm và cách xử lý trong đề tài}
\label{tab:query-modes}
\begin{tabularx}{\textwidth}{|>{\raggedright\arraybackslash}p{2.6cm}|L|L|L|}
\hline
\textbf{Hình thức} & \textbf{Đầu vào} & \textbf{Ưu điểm và thách thức} & \textbf{Cách xử lý trong đề tài} \\
\hline
Ảnh tham chiếu (ReID) & Ảnh đã cắt chứa một người & Giàu thông tin thị giác; nhạy với góc nhìn, ánh sáng, che khuất & Mã hóa bằng bộ mã hóa ảnh \\
\hline
Mô tả văn bản (TBPS) & Câu mô tả tiếng Anh & Không cần ảnh; mô tả có thể mơ hồ, thiếu chi tiết & Mã hóa bằng bộ mã hóa văn bản \\
\hline
Thuộc tính & Giới tính, loại và màu trang phục trên, dưới, vật mang theo & Dễ nhập, rõ ràng; chỉ diễn đạt được các thuộc tính có sẵn & Chuyển thành câu tiếng Anh có cấu trúc, rồi xử lý như mô tả văn bản \\
\hline
\end{tabularx}
\end{table}

\begin{itemize}
    \item \textbf{Một không gian biểu diễn chung.} RaSa~\cite{bai2023rasa} mã hóa ảnh và văn bản vào cùng một không gian vector, nên ảnh người, ảnh truy vấn, câu mô tả và câu từ thuộc tính so sánh trực tiếp được (Mục~\ref{sec:3.5}).
    \item \textbf{Đơn vị kết quả là một lần xuất hiện.} Người được phát hiện và theo vết (Mục~\ref{sec:3.2} và~\ref{sec:3.3}); mỗi lần xuất hiện liên tục trên một camera (track) có một khung hình đại diện và một vector, giảm số phần tử cần so khớp và tránh kết quả trùng.
    \item \textbf{Hai giai đoạn tách biệt.} Lập chỉ mục chạy nền và lưu đặc trưng; truy vấn chỉ mã hóa mô tả và tìm các vector gần nhất (Mục~\ref{sec:3.6}).
    \item \textbf{Con người ra quyết định cuối cùng.} Hệ thống chỉ xếp hạng; người dùng đánh giá từng kết quả trên hình ảnh gốc trước khi lưu.
\end{itemize}
